\def\ChapterCode{NEU}
\chapter
[\CO{[NEU]} Neurologische Störungen]
{Neurologische Störungen}
\ChapterIntro
{%
Neurologische Störungen sind Störungen bzw.
Erkrankungen und Notfälle die das Nervensystem betreffen. Sie
gehen oft mit einer Störung des Bewusstseins oder anderer
Funktionen des Nervensystems einher. Die \EE{Ursachen} von
neurologischen Erkrankungen sind sehr vielfältig. }
{%
\begin{description}
  \item[Maintainer:] Sebastian Gabriel
  \item[Autoren:] Diverse
  \item[Reviewer:] Standard-Reviewprozess
  \item[Version:] {\VarDocVersionTypeName} ({\VarDocVersionTypeDescription})
  \item[SHA1:] {\DocGitCommit}
\end{description}

{\TxTeilkapitelAASS}
}

\bigskip

\noindent Die folgende Übersicht nennt einige
Beispiele für Ursachen neurologischer Störungen, ohne
Anspruch auf Vollständigkeit:
\begin{itemize}
  \item \E{Schlaganfall}
  \item \E{Vergiftungen} können fast alles verursachen, inkl.
  neurologischer
  Symptome.
  \item \E{Schädel-Hirn-Trauma} (SHT); Ein Schlag auf
  den Kopf erhöht \emph{nicht} das Denkvermögen \ldots
  \item \E{Stoffwechselstörungen} welche Einfluss auf das
  Nervensystem nehmen, z.\,B. Zuckerentgleisungen
  \item \E{Infektionen} können das Nervensystem beeinflussen,
  z.\,B. Verwirrtheit bei Fieber
  \item \E{Alter} beeinflusst auch das Nervensystem, im Rahmen
  einer Demenz kommt es z.\,B. zum Abbau der Denkleistung
  \item \E{Krampfanfälle} können vom zentralen Nervensystem ausgehen
  \item \E{Erkrankungen des Stützapparates} Aufgrund der
  räumlichen Nähe zum Stütz- und Bewegungsapparates kann das
  Nervensystem in Mitleidenschaft gezogen werden, z.\,B. bei
  Bandscheibenvorfällen
\end{itemize}

\clearpage
\Text{{\TxQuerverweise}}{%
\begin{itemize}
  \item Infektiöse (bakterielle) Meningitis: \dotfill
  \FullPrettyRef{topic:meningitis}
\end{itemize}
}{}{}{}{}




\section{Schlaganfall, Apoplektischer Insult, Hirninfarkt}
\index{Schlaganfall}
\index{Insult}
\index{Apoplexie}
\index{Apoplektischer Insult}
\label{topic:insult}
\label{topic:schlaganfall}

\HeadingSub{\Lang{Syn.} Apoplektischer Insult, Apoplexie}

% \Definition{Siehe Arten.}

\TextSlide{Arten}
{%
Grundsätzlich gibt es 2 Formen, die aber vor Ort nicht
unterschieden werden können:
\begin{itemize}
  \item \EE{\T{Trockener Insult}}:
  \I[Ischämie!trockener Insult]{Ischämie}. Verschluss eines
  Hirngefäßes mit anschließendem Absterben der versorgten Hirnregion
  \item \EE{\T{Blutiger Insult}}: \I{Hirnblutung}, z.\,B.
  durch Platzen eines Aneurysmas eines Hirngefäßes, kann spontan
  oder aufgrund eines Traumas auftreten.
\end{itemize}

}
{
\begin{itemize}
  \item \enquote{Trockener} Insult
  \item \enquote{Blutiger} Insult
\end{itemize}

}{}{}{}

% TODO Anamnese Insult vor Stroke-Anmeldung

\TextSlide{TIA}
{%
\index{TIA}
\index{Transitorisch-Ischämische Attacke}
\label{topic:tia}
\T{TIA} ist die Abkürzung für \I{Transitorisch-Ischämische Attacke}.
Von einer TIA spricht man, wenn die Symptome  innerhalb von
24 Stunden vergehen. Solange noch Symptome bestehen, kann
man nicht von einer TIA sprechen, der Patient ist wie bei
jedem anderen Schlaganfall zu behandeln.
\bigskip
}
{
\begin{itemize}
  \item Transitorisch-Ischämische Attacke.
  \item Symptome <\,24\Hour*.
\end{itemize}

}{}{}{}


\TextSlide{\TxSymptome}
{%
\label{topic:hirndruckzeichen}
Die diversen \EE{Funktionsareale}
(Sprachzentren, Sehzentrum, motorische Zentren, \ldots)
sind in verschiedenen Regionen des Hirns verteilt. Je
nachdem welches Gebiet durch die Ischämie oder Blutung
betroffen ist, können die Symptome sehr unterschiedlich
sein.  Typische Störungen sind:

\begin{itemize}
  \item \EE{Motorik}: Störungen der Kraft und Beweglichkeit
  sind sehr häufig.
  Charakteristisch für einen
  Schlaganfall sind sog. \T{Halbseitenzeichen}: Es kommt
  dabei zu einer \E{einseitigen Kraftminderung}, dies kann
  bis zu einer \E{schlaffen, kompletten Lähmung} führen. Ein
  \E{herabhängender Mundwinkel} ist ein Zeichen einer
  Lähmung der
  Gesichtsmuskulatur. Durch die Lähmungen kann es zu
  \EE{Schluckstörungen} kommen, es besteht dann eine erhöhte
  Aspirationsgefahr.
  
  Der \T{Herdblick} ist eine Blickdiagnose: Es kann zu
  Blicklähmungen kommen, der Patient blickt in Richtung des
  \enquote*{Herdes des Geschehens}.
  
  \item \EE{Sprache}: Störungen der Sprache sind ebenfalls
  sehr häufig. Oft nimmt man eine \EE{verwaschene Sprache}
  wahr, welche (fälschlich) an Trunkenheit erinnert.
  \Z{Weiters kann es zu einer \enquote{wirren Sprache}
  kommen: Worte fehlen dabei oder Silben werden vertauscht.
  Im Extremfall kann die Sprache gänzlich \enquote{verloren}
  gehen. Auf der anderen Seite kann die Sprache zwar
  erhalten, aber das Sprach\E{verständnis} geschädigt sein.}
  
  \item \EE{Wahrnehmung}:
  Es kann zu Wahrnehmungsstörungen wie z.\,B.
  \EE{Sehstörungen}, plötzliche Erblindung, Doppelbilder,
  Gesichtsfeldausfälle kommen.
  
  \item \EE{Bewusstsein}: Störungen des Bewusstseins sind
  gefährlich und müssen bei Insult-Patienten immer erwartet
  werden. Der Bewusstseinszustand kann sich auch plötzlich
  verschlechtern.
  
  \item \EE{Allgemein}: Weiters können unspezifische
  Symptome wie z.\,B. Schwindel und Kopfschmerzen auftreten. Ein
  plötzlicher, \EE{donnerschlagartiger Kopfschmerz} ist ein
  klassisches Symptom für eine spontane Hirnblutung.
  
  \item \EE{\T{Hirndruckzeichen}}: Bei einer Erhöhung des
  Hirndrucks, z.\,B. aufgrund einer Hirnblutung, kommt es zu
  einer typischen Kombination von Symptomen: 
  
  \parpic[4cm][r]{\includegraphics[width=4cm]{./images/trauma/pupillendifferenz.jpg}}
  Meist das erste Symptom sind \EE{Kopfschmerzen}, gefolgt
  von Übelkeit. In weiterer Folge treten
  \EE{Bewusstseinsstörungen} auf. Der Patient ist
  \EE{hyperton, aber bradykard} (\EE{RR\,\Sign{↑},
  HF\,\Sign{↓}}), der Puls ist langsam, aber äußerst stark
  spürbar (\I{Druckpuls}). Die \EE{Atmung} kann \EE{unregelmäßig} sein. Weiters
  ist die \EE{Ungleichheit der Pupillen} \Z{(Anisokorie)\index{Anisokorie|Z}}, bzw.
  eine verlangsamte Lichtreaktion
  typisch.\ZFootnote{\E{Erhöhter Hirndruck}:In schweren
  Fällen kommt es zu sog. Strecksynergismen, d.\,h.  der
  Patient reagiert bei Schmerzreiz nur mit einem
  ungerichtetem Strecken der Extremitäten.}
  \nocite{StriebelAnaesthesie2} %[Bd.2, S.1344]
  
  \Z{Hirndruckzeichen werden nicht nur durch
  Hirnblutungen verursacht, auch z.\,B. beim \Z{Hirnödem}
  oder Schädel-Hirn-Trauma (\Abk{SHT}) können sie
  vorliegen.}
  \item \EE{Blutdruck}: Der Blutdruck ist in der
      Akutsituation in der Regel erhöht und sinkt spontan
      innerhalb der nachfolgenden Stunden. Es wird allgemein
      empfohlen Blutdruckwerte bis \RRmmHg{220}{120} zu
      tolerieren, wenn nicht cardiopulmonale Gründe (Myocardischämie,
      Herzinsuffizienz, \ldots) andere Grenzwerte erfordern.
\end{itemize}

Präklinisch kann man nicht sicher zwischen einer Ischämie
(trockener Insult) und einer Blutung unterscheiden.
Auf jeden Fall ist die \EE{Bewusstseinslage} sehr engmaschig
zu überwachen; eine rasche Verschlechterung ist jederzeit
möglich! 
}
{
Je nach betroffener Hirnregion sehr unterschiedlich.
\begin{itemize}
  \item Halbseitenzeichen (Schwäche, Lähmung)
  \item Herdblick
  \item Verwaschene Sprache
  \item Sehstörungen
  \item Bewusstseinsstörungen
  \item Kopfschmerzen
  \item Hirndruckzeichen
  \begin{itemize}
    \item Bewusstseinsstörungen
    \item RR\,\Sign{↑} \& HF\,\Sign{↓}
    \item Ungleichweite Pupillen, verlangsamte Lichtreaktion
    \item Strecksynergismen
  \end{itemize}
  
\end{itemize}
}{}{}{}


\TextSlide{Differentialdiagnose}
{
Die wichtigste Differentialdiagnose ist eine Störung des
\EE{Blutzucker}-Stoffwechsels. Sowohl die Hyper-, als auch
die Hypoglykämie können auch fokalneurologische Symptome verursachen
  und einen Schlaganfall vortäuschen, daher
muss der Blutzucker immer gemessen werden!
Ebenso ist immer
an \EE{Vergiftungen} zu denken.

Die Abgrenzung zu anderen neurologischen Erkrankungen ist
ohne weitergehende Untersuchungen oft schwierig bzw. schwer
möglich. 
}{
\begin{itemize}
  \item Blutzucker ${\downarrow}$, ${\uparrow}$
  \item Sonstige neurologische Erkrankung.
  \item Vergiftungen
\end{itemize}
}

\TextSlide{Spezielle Anamnese}
{%
Vor Anmeldung zu einer Stroke-Unit müssen folgende
Informationen erhoben worden sein:

\begin{itemize}
  \item Zeitpunkt des Auftretens der Symptome
  \item 
\end{itemize}

}{
\begin{itemize}
  \item 
\end{itemize}

}{}{}{}

% M %%%%%%%%%%%%%%%%%%%%%%%%%%%%%%%%%%%%%%%%%%%%%%%%%%% M %
% Insult
\ProcedurePrintDefault{NI64XX0C}


\section{Vom Gehirn ausgehende Krämpfe -- Zerebrale Krampfanfälle}
\label{topic:zerebrale-krampfanfaelle}
\label{topic:krampfanfall}




\TextSlide{{\TxBeschreibung} und Einteilung}
{%
\DefinitionInPara{Bei vom Gehirn ausgehenden motorischen Krampfanfällen
kommt es zur unkontrollierten Entladung von Nervenzellen im Gehirn zu Krämpfen der 
Skelettmuskulatur. Mögliche Ursachen sind akute oder chronische Erkrankungen des
Gehirns.}
Die zerebralen Krampfanfälle kann man auf verschiedene Arten
einteilen:
\begin{enumerate}
  \item Nach Art:
  \begin{itemize}
    \item \EE{Tonisch}: \enquote{Verkrampfung}, erhöhter
    Muskeltonus
    \item \EE{Klonisch}: Zuckungen
    \item \EE{Tonisch-klonisch}: Kombination aus beidem, das
    klassische Erscheinungsbild eines Krampfanfalles
  \end{itemize}
  \item Nach Lokalisation:
  \begin{itemize}
    \item \EE{Fokaler Anfall}: Der Krampfanfall ist auf
    eine Körperregion beschränkt, siehe
    \ref{topic:krampfanfall-fokaler}
    \item \EE{Generalisierter Anfall}: Er betrifft den
    ganzen Körper, siehe
    \ref{topic:krampfanfall-generalisierter}
    
  \end{itemize}
\end{enumerate}
}{%
\begin{itemize}
  \item Nach Art:
  \begin{description}
    \item Tonisch
    \item Klonisch
    \item Tonisch-klonisch
  \end{description}
  \item Nach Lokalisation:
  \begin{description}
    \item Fokaler Anfall
    \item Generalisierter Anfall
  \end{description}
\end{itemize}
}{}{}{}

\Text{Ursachen}
{%
Für zerebrale Anfälle gibt es eine Reihe von
unterschiedlichen Ursachen die direkt oder indirekt das
Gehirn betreffen. Einige betreffen die Hirn- und
Nervenstruktur, andere sorgen dafür, dass das Hirn zu wenig
Sauerstoff und andere Nährstoffe erhält. Die häufigsten
Ursachen und Grunderkrankungen sind im Folgenden aufgeführt:

\begin{itemize}
  \item \EE{Epilepsie}: Die Epilepsie ist eine chronische
  Erkrankung welche durch wiederkehrende Episoden von
  zerebralen Krampfanfällen gekennzeichnet ist.
  \item \EE{Vergiftungen}: (Medikamente/Drogen/...) können
  wieder fast alles machen, also auch Krampfanfälle. \citep{Supady2009}
  \item \EE{Hirntumor}: Durch den Tumor wird das Hirn
  beschädigt, auch dadurch kann es zu Anfällen kommen.
  \item \EE{Narbengewebe}: nach Insult oder Hirnverletzung.
  \item \EE{Hirnblutung}
  \item \EE{Sauerstoffmangel}: z.\,B. Kollaps, Ischämie
  aufgrund eines Gefäßverschlusses
  \item \EE{Zuckerstoffwechselstörung}: Der Dauerbrenner.
  Eine Blutzuckerentgleisung kann fast alle neurologischen Symptome vortäuschen!
  \item \EE{angeboren}: ohne organisch fassbare  Ursache
  \item \EE{Fieberkrampf}: Fieber bei Kleinkindern
  \item \EE{Infektionen}
  \item \EE{Entzug}: Entzug von Alkohol oder anderen Drogen
\end{itemize}
}{
}{}{}{}




\Text
{Fokaler Anfall}
{%
\index{Krampfanfall!fokaler}
\label{topic:krampfanfall-fokaler}
Der Krampfanfall ist auf eine Körperregion beschränkt,
z.\,B. auf eine Hand. Der Patient kann bei Bewusstsein sein.
Ein fokaler Anfall kann sich ausweiten und in einen
generalisierten Anfall übergehen.

\Z{Der fokale Anfall wird
auch \enquote{petit mal}\ZFootnote{Petit mal: \emph{franz.:} \enquote{kleines
Unglück}} genannt.}}
{}
{}
{}
{}

\TextSlide
{Generalisierter Anfall}
{%
\index{Krampfanfall!generalisierter}
\label{topic:krampfanfall-generalisierter}
Der generalisierte Krampfanfall \Z{(\enquote*{grand mal}\ZFootnote{Grand mal: \emph{franz.:} \enquote{großes
Unglück}})} betrifft den ganzen Körper. Häufig kommt es während des Anfalles zu einem Zungenbiss,
außerdem zum Einnässen oder Einkoten. 
Der generalisierte Anfall läuft typischerweise in mehreren Phasen ab:

\subparagraph{Beginn}
Dem Anfall kann ein
\E{Dämmerzustand} 
vorausgehen, \Z{eine sog. \enquote{Absence}}. Der Patient wirkt
dabei plötzlich abwesend und \emph{\enquote{irgendwie komisch}}.
Dann verliert er das Bewusstsein und die Muskeln werden
anfangs schlaff. Wenn der Patient gestanden ist, fällt er
ungebremst auf den Boden (Cave: \EE{Verletzungen!}).
Manchmal kommt es vor, dass der Patient einen
\I{Initialschrei} von sich gibt, der
gruselig klingen kann. 

\subparagraph{Krampf}
Dann kommt es zu
\EE{tonisch-klonischen Krämpfen},
die den ganzen Körper betreffen. Während des Anfalls ist der
Patient nicht bei Bewusstsein, es besteht außerdem i.\,d.\,R.
ein Atemstillstand. Er kann jedoch starke Kräfte entwickeln, es
besteht massive
Verletzungsgefahr.

\subparagraph{Danach}
Der Krampf dauert meist wenige Minuten und hört von selbst
auf. Der Patient ist danach erstmal bewusstlos, im weiteren
Verlauf mit mehr oder weniger großem Aufwand erweckbar. Er
befindet sich dann in der sogenannten
\T{Nachschlafphase}\index{Nachschlafphase}.
Der Patient \EE{klart im Verlauf mehr und mehr auf}. Es ist
wichtig zu versuchen mit dem Patienten in Kontakt zu treten
um seinen Wachheitsgrad und dessen
\EE{Verlauf} beurteilen zu können. Für den Anfall selbst besteht
Amnesie! Daher muss man den Patienten aufklären was passiert
ist, wo er sich befindet, etc. Ein \EE{Neuro-} und
\EE{Traumacheck} ist nach \underline{\E{jedem}} Anfall
unbedingt notwendig! 
}
{
\begin{itemize}
  \item Betrifft ganzen Körper
  \item Zungenbiss
  \item Einnässen, Einkoten
  \item Beginn:
  \begin{compactitem}
    \item Dämmerzustand (Absence)
    \item Initialschrei
    \item Evtl. Sturz und Folgeverletzungen
  \end{compactitem}
  \item Krampf:
  \begin{compactitem}
    \item Tonisch-klonische Krämpfe
    \item Verletzungsgefahr
  \end{compactitem}
  \item Danach:
  \begin{compactitem}
    \item \EE{Nachschlafphase}
    \item \EE{Aufklaren}
    \item Amnesie
    \item Neurocheck
    \item \EE{Traumacheck} nach \underline{\E{jedem}} Krampfanfall
  \end{compactitem}
\end{itemize}

}
{}
{}
{}


\TextSlide{Status Epilepticus}
{\index{Status!epilepticus}
Der Status Epilepticus ist eine
lebensgefährliche Steigerung
des generalisierten Krampfanfalls. Er besteht dann, wenn der
Anfall nicht nach kurzer Zeit vergeht oder innerhalb kurzer
Abstände mehrere Anfälle auftreten (z.\,B. 2 oder mehr Anfälle
innerhalb von 24 Stunden). Der Krampf muss dann mit Medikamenten
durchbrochen werden.
}
{
\begin{itemize}
  \item Lange Dauer
  \item Wiederholtes Auftreten ($\geq$\,2\,/\,24\,h)
  \item Akute \Ca{Lebensgefahr!}
\end{itemize}
}{}{}{}

\VARIANT{VAR-RETTUNGSDIENST}
{\Fazit{Im Rettungsdienst gilt praktisch jeder Krampfanfall,
den das Personal selber miterlebt, als Status-Verdächtig.}}



% M %%%%%%%%%%%%%%%%%%%%%%%%%%%%%%%%%%%%%%%%%%%%%%%%%%% M %
% Krampfender Patient
\ProcedurePrintDefault{NR56081C}


% M %%%%%%%%%%%%%%%%%%%%%%%%%%%%%%%%%%%%%%%%%%%%%%%%%%% M %
% Nicht-(mehr) krampfender Patient
\ProcedurePrintDefault{NR56082C}


\subsection{Besondere Krampfanfälle}

\Layout{60}{Fieberkrampf}
{%
\label{topic:fieberkrampf}%
\index{Fieberkrampf}%
\index{Krampfanfall!Fieberkrampf}%
%
\EE{Kinder} unter 4~Jahren können bei schnellem, hohen
Fieberanstieg ({\textgreater}\,39{\DegC}) 
\E{Fieberkrämpfe} entwickeln. Sie entsprechen einem vom
Gehirn ausgehenden Krampfanfall und sind grundsätzlich genau
so zu behandeln.

Mitunter ist die Neigung des Kindes zu Fieberkrämpfen schon
bekannt und es wurden bereits spezielle Zäpfchen für den
Bedarfsfall verschrieben (und bei Eintreffen eventuell schon
von den Eltern gegeben).
}{
\begin{itemize}
  \item Kinder {\textless} 4 Jahre
  \item Hohes Fieber
  \item Wie generalisierter Anfall
\end{itemize}
}
{./images/gabriel-sebastian-aass/fieberthermometer.jpg}
{}
{GABS}



\TextSlide{Eklamptischer Krampfanfall während der Schwangerschaft}
{%
In Folge einer Präeklampsie (EPH-Gestose)
kann es zu einem sog. \E{eklamptischen Krampfanfall}
kommen. Dieser wird unter \FullPrettyRef{topic:eklampsie}
besprochen.
}{
\begin{itemize}
    \item Folge einer Präemklampsie.
    \item \FullPrettyRef{topic:eklampsie}.
\end{itemize}
}{}{}{}





% \clearpage % TEMP temp clearpage
\section{Lumbago, Lumbo-ischialgie und Bandscheibenvorfall} 
\HeadingSub{\Z{(Discusprolaps)}} 
\label{topic:lumbago}

\TextSlide{\TxSymptome}{%
Verschiedene Ursachen wie z.\,B. eine Nervenreizung, ein
Bandscheibenvorfall (siehe unten), Erkrankungen der Wirbelsäule oder ein Unfall
können zu einem Druck auf den \E{Ischias-Nerv} führen, 
welcher starke Schmerzen im Bereich der Lendenwirbelsäule 
(\Abk{LWS}) verursachen kann. Die Schmerzen strahlen oft ins Gesäß und in
das betroffene Bein bis zum Fußaußenrand aus und werden von
Gefühlsstörungen in diesem Gebiet bzw. von einer Verkrampfung der
Rückenmuskulatur begleitet. Diese Symtomatik nennt man \I{Ischialgie} bzw. \I{Lumbalgie}. 
Bei plötzlichem Auftreten
spricht man von Lumboischialgie (\T{Lumbago}) oder umgangssprachlich vom
{\I{Hexenschuss}}.

Bei vorgeschädigten Bandscheiben (z.\,B. durch jahrelange
Fehlhaltung oder durch falsches Heben großer Lasten) kann es
zu einer (plötzlichen) Verschiebung des Bandscheibenkerns in Richtung Rückenmark und somit zu einem Druck auf dieses
kommen. Dieser Druck verursacht ebenfalls Schmerzen,
Sensibilitätsstörungen oder sogar Lähmungen. Tritt dieses Ereignis plötzlich auf, spricht man von einem
\T{akuten Bandscheibenvorfall}. Häufig ist die Lendenwirbelsäule (LWS-Bereich)
betroffen.
}{
\begin{itemize}
  \item Schmerzen in Lendengegend
  \item Evtl. Austrahlung in Bein
  \item Gefühlsstörungen
\end{itemize}
}{}{}{}

% \noindent Unter Ischialgie (\T{Lumbalgie}, Lumboischialgie,
% \emph{\enquote{\T{Lumbago}}}) versteht man Schmerzen im
% Versorgungsbereich des Ischias-Nerv, welche u. a. durch eine
% Entzündung oder durch Reizung bzw. Kompression des Nervs
% oder seiner Wurzeln verursacht werden können.
% 
% \TextSlide{Ursachen}{%
% können eine Reizung bzw. Druck
% im Bereich der Lendenwirbel bzw. der Kreuzbeinwirbel,
% ein Bandscheibenvorfall, sowie andere Erkrankungen der Wirbelsäule, aber auch
% ein Unfall sein. 
% 
% \bigskip
% 
% \bigskip
% 
% }{
% \begin{compactitem}
%   \item Nervenreizung
%   \item Bandscheibenvorfall
%   \item Erkrankungen der Wirbelsäule
%   \item Unfall
% \end{compactitem}
% }{}{}{}
% 
% \TextSlide{\TxSymptome}{%
% Das Leitsymptom sind Schmerzen in der Lendengegend, die in
% das betroffene Bein bis zum Fußaußenrand ausstrahlen,
% Schonhaltung, lokale Druck- und Klopfempfindlichkeit der
% Wirbelsäule, Sensibilitätsstörungen und eventuell Lähmung
% der Zehenmuskulatur. Je nach Schwere und zeitlicher
% Entwicklung (langsames/plötzliches Auftreten) bezeichnet man
% das Zustandsbild als Lumboischialgie oder Lumbago
% (Hexenschuss). Gegebenenfalls ist im Krankenhaus ein
% Bandscheibenvorfall abzuklären.
% }{
% \begin{itemize}
%   \item Schmerzen in Lendengegend
%   \item Evtl. Austrahlung in Bein
%   \item Gefühlsstörungen
% \end{itemize}
% }{}{}{}



\ChapterEnd